\chapter*{Preface}
\addcontentsline{toc}{chapter}{Preface}
\markboth{Preface}{Preface}

A fluctuating quantity can disappear from a description while its effects
remain. Consider the displacement $x$ of a spring in thermal equilibrium.
Choose units in which its energy is $x^2/2$, and write $h>0$ for the thermal
energy. The probability density of the displacement is proportional to
$e^{-x^2/(2h)}$. A polynomial $f(x)$ represents a measurement, or
\emph{observable}. Its average is
\[
 \langle f\rangle_h=
 \frac{\displaystyle\int_{\mathbb R}f(x)e^{-x^2/(2h)}\,dx}
      {\displaystyle\int_{\mathbb R}e^{-x^2/(2h)}\,dx}.
\]
The denominator is positive and finite. All polynomial moments exist:
outside a sufficiently large interval, any fixed power of $|x|$ times
$e^{-x^2/(2h)}$ is bounded by a constant times $e^{-|x|}$.
Symmetry gives $\langle x\rangle_h=0$. Yet the spring moves, and
$\langle x^2\rangle_h=h$.
Replacing the displacement by its mean therefore loses a measurable effect.

We want to simplify a system while retaining its operations on observables.
Measurements can be multiplied, moved, or brought together.
Symmetries can act before or after elimination.
The equalities comparing these procedures contain more information than
the separate answers. For the spring, one integration by parts supplies them.

If $p$ is a real polynomial, then $p(x)e^{-x^2/(2h)}$ tends to zero at
both ends of the real line. Integrating its derivative gives
\[
 \boxed{\quad\langle xp\rangle_h=h\langle p'\rangle_h.\quad}
\]
Taking $p=1,x,x^2,\ldots$ determines every moment. The odd moments vanish,
and, for $n\ge1$, the even moments satisfy
\[
 \langle x^{2n}\rangle_h
 =(2n-1)h\langle x^{2n-2}\rangle_h
 =1\cdot3\cdots(2n-1)h^n,
 \qquad \langle1\rangle_h=1.
\]
Thus the same equation that eliminates $x$ also determines its surviving
fluctuations.

We can retain the equation itself. Set
$B_h p=xp-hp'$. Every $B_h p$ has average zero, and every polynomial of
average zero has this form. To see the converse, remove the leading term
$a x^m$ of a nonconstant polynomial by subtracting $aB_h(x^{m-1})$.
The degree decreases. Repetition writes any polynomial uniquely as
$c+B_h p$. The constant is its average, and uniqueness follows because
$B_h p$ has degree $\deg p+1$ whenever $p\ne0$.
The polynomial $p$ records how the relation was obtained.

Multiplying that relation by another polynomial changes it:
\[
 \boxed{\quad (B_h p)q=B_h(pq)+hpq'.\quad}
\]
This follows by expanding $(pq)'=p'q+pq'$. In particular,
$(B_h1)x=B_hx+h$. The zero-average polynomial $x=B_h1$ becomes
the polynomial $x^2$, whose average is $h$.
Consequently, polynomials with the same average cannot be multiplied
as equivalence classes using their ordinary product.
The correction $hpq'$ is part of the elimination rule.

At $h=0$, the rule becomes ordinary division by $x$:
$f(x)=f(0)+xp(x)$. The terms involving $h$ describe exactly how a
relation at zero fluctuation changes when fluctuations are admitted.
This is a useful reason to retain a relation together with its primitive,
the element mapped to it by $B_h$.
In more general problems there are several possible primitives.
Their differences, their products, and the equations comparing those
products then become mathematical data in their own right.

\medskip

Elimination becomes richer when some variables remain. Let $y$ be a
retained displacement and let $x$ be an independent Gaussian displacement
of variance $h$. Averaging a polynomial of their sum defines
\[
 T_h f(y)=\langle f(y+x)\rangle_h
 =\sum_{k\ge0}\frac{h^k}{2^k k!}f^{(2k)}(y).
\]
Taylor's formula and the moments above give the equality. The sum is finite
for every polynomial. The operator retains the mean together with its
fluctuation corrections.

The same finite formula defines $T_h$ for any real $h$, including negative
values. Negative $h$ here specifies a polynomial operator, with no
probability interpretation. Composing the formulas gives
\[
 T_aT_b=T_{a+b},\qquad T_0=1,\qquad T_h^{-1}=T_{-h}.
\]
Multiply the coefficient of $f^{(2m)}$ by $2^m m!$.
It becomes the binomial sum
$\sum_{j=0}^m\binom mj a^jb^{m-j}=(a+b)^m$.
For $a,b>0$, eliminating independent Gaussian displacements adds
their variances.

Because $T_h$ is invertible on polynomials, it carries ordinary
multiplication to a multiplication on the averaged descriptions:
\[
 f\star_h g=T_h\bigl((T_{-h}f)(T_{-h}g)\bigr).
\]
Thus $T_h(fg)=(T_hf)\star_h(T_hg)$. Explicitly,
\[
 \boxed{\quad
 f\star_h g=\sum_{k\ge0}\frac{h^k}{k!}f^{(k)}g^{(k)}.
 \quad}
\]
Here is a direct verification. Temporarily give the two factors variables
$u,v$, and restrict to $u=v=y$ after differentiating.
Differentiation in $y$ then acts as $\partial_u+\partial_v$.
Write $T_h=\exp(h\partial_y^2/2)$, where the exponential means its
terminating power series on polynomials. The three commuting operators
in the definition of $\star_h$ combine to give
\[
 \exp\!\left(\frac h2(\partial_u+\partial_v)^2
                 -\frac h2\partial_u^2-\frac h2\partial_v^2\right)
 =\exp(h\partial_u\partial_v).
\]
Expanding this last exponential proves the formula.

Each derivative in one factor is paired with a derivative in the other.
For example,
\[
 y\star_h y=y^2+h,\qquad
 (y\star_h y)\star_h y=y^3+3hy
 =y\star_h(y\star_h y).
\]
Associativity holds for all polynomials because applying $T_{-h}$ turns
either parenthesization into an ordinary product of three polynomials.
For $r$ factors the same argument gives
\[
 f_1\star_h\cdots\star_h f_r
 =\left.
 \exp\!\left(h\sum_{i<j}\partial_{y_i}\partial_{y_j}\right)
 \prod_{i=1}^r f_i(y_i)\right|_{y_1=\cdots=y_r=y}.
\]
All sums still terminate. For monomials, the derivatives count choices of
paired factors and the factorials remove repeated orderings of those pairs.
Thus a single two-factor correction already
has consequences for every number of measurements.

There is an elementary way to retain all these composition equations.
For an integer $r\ge1$, let $W_r$ be the real vector space of finite
linear combinations of words $[f_1|\cdots|f_r]$ of $r$ real polynomials,
imposing linearity in every letter. Let $W_0$ be the one-dimensional
space spanned by the empty word $[\,]$. An element of
$W=\bigoplus_{r\ge0}W_r$ is a finite sum of elements of these spaces.
Assign degree $-r$ to $W_r$.

Join adjacent letters using $\star_h$, taking alternating signs from
left to right. Denote this linear map by $b:W\to W$, and set it to
zero on words of lengths zero and one. It raises degree by one.
Its first nonzero formulas are
\[
 \begin{aligned}
 b[f|g]&=[f\star_h g],\\
 b[f|g|k]&=[f\star_h g|k]-[f|g\star_h k].
 \end{aligned}
\]
Applying $b$ again to the second line gives
$[(f\star_h g)\star_h k-f\star_h(g\star_h k)]=0$.
On longer words, two disjoint joins occur twice with opposite signs.
Overlapping joins cancel by the same three-letter calculation.
Hence $b^2=0$ at every word length.

This is the multiplication part of the \emph{bar construction}.
The degree-one map $b$, whose square is zero, is called a differential.
The Gaussian product makes the three-letter defect vanish exactly.
A weaker comparison would retain an explicit correction to such a defect.
For a bilinear multiplication $\mu:V\times V\to V$ on a real vector
space $V$, choose another real vector space $E$ and a linear map
$L:E\to V$. A trilinear correction $K:V\times V\times V\to E$
would have to satisfy
\[
 \mu(\mu(a,b),c)-\mu(a,\mu(b,c))=L\bigl(K(a,b,c)\bigr).
\]
Here $K(a,b,c)$ is a primitive of the associator under the specified map
$L$. This equation is a requirement on $K$, and supplies no existence
claim. It also leaves the comparison of four or more inputs unsettled.
The search for corrections compatible with every grouping motivates
higher associativity.

There is also an elementary motivation for a dual construction.
A linear functional on $W$ is a linear map $W\to\mathbb R$.
Write $W^\vee$ for the space of all these functionals.
Cutting a word at each possible position defines their product:
\[
 (\lambda*\eta)[f_1|\cdots|f_r]
 =\sum_{i=0}^r
   \lambda[f_1|\cdots|f_i]\,
   \eta[f_{i+1}|\cdots|f_r].
\]
The term $i=0$ uses the empty left word, and $i=r$ uses the empty
right word. Extend by linearity to $W$.
Every evaluation is a finite sum and respects linearity in each letter.
Both $(\lambda*\eta)*\theta$ and $\lambda*(\eta*\theta)$ evaluate
on a word as the sum
\[
 \sum_{0\le i\le j\le r}
 \lambda[f_1|\cdots|f_i]\,
 \eta[f_{i+1}|\cdots|f_j]\,
 \theta[f_{j+1}|\cdots|f_r].
\]
Thus the product is associative. Its unit takes value one on the empty
word and zero on all positive lengths.

The transpose of $b$ means precomposition:
$b^\vee\lambda=\lambda\circ b$. It is a linear map $W^\vee\to W^\vee$,
and $(b^\vee)^2\lambda=\lambda\circ b^2=0$.
On two letters it gives
$(b^\vee\lambda)[f|g]=\lambda[f\star_h g]$.
Thus a coefficient measured after joining becomes a two-input coefficient,
while cutting has supplied a product of coefficient functionals.
These elementary operations motivate the exchange between joining and
splitting in bar and cobar duality.
They use finite words throughout. Admitting infinite sums of words
requires further choices about which coefficient evaluations are defined.

\medskip

Even counting the polynomials forces this issue. Give independent variables
$x_1,x_2,\ldots$ positive integer weights $d_1,d_2,\ldots$.
A monomial $\prod_j x_j^{n_j}$ has weight $\sum_j d_jn_j$, with only
finitely many $n_j$ nonzero. If only finitely many variables have weight
below any prescribed bound, then the number of monomials of each weight
is finite. Choosing their exponents gives
\[
 \sum_{m\ge0}(\text{number of monomials of weight }m)t^m
 =\prod_j(1+t^{d_j}+t^{2d_j}+\cdots)
 =\prod_j\frac1{1-t^{d_j}}.
\]
Each coefficient uses finitely many choices. This makes the infinite
product a well-defined formal series, without assigning a numerical
value to $t$.

The product counts available polynomial measurements. It contains no
choice of Gaussian average. Two variables can have independent
fluctuations, or they can be correlated, with the same monomial count.
For example, take independent centered variables $x,u$ of variance $h$
and set $v=cx+u$. The polynomials in $x,v$ have the same basis count
for every real $c$, whereas $\langle xv\rangle=ch$.
The coefficient needed to combine two measurements has disappeared
from the counting series.

A formal series also retains more than any one numerical substitution.
Consider $f(y)=\sum_{n\ge0}y^{2n}$. The constant coefficient of
$T_hf$, denoted by $[y^0]T_hf$, is formally
\[
 [y^0]T_hf=\sum_{n\ge0}1\cdot3\cdots(2n-1)h^n,
\]
where the term for $n=0$ is $1$. For fixed powers of $h$ and $y$,
only one term of $f$ contributes, so this calculation is well-defined
as a series in two formal variables. For every numerical $h>0$, the
displayed terms eventually increase: the ratio of consecutive terms is
$(2n+1)h$. The numerical series diverges.
Agreement of all finite coefficient truncations therefore does not
supply a convergent average at a fixed positive variance.

Completing a space means specifying which infinite coefficient lists
are admitted, and which finite portions determine their approximations.
For a nonnegative integer $N$, reduction modulo $h^{N+1}$ retains
only the first $N+1$ powers
of $h$. Results at successive orders must agree under these reductions.
An operation on the completion must have a finite prescription for each
retained coefficient. Substituting a number for $h$ is a further operation,
with its own convergence question.
To give the Gaussian variables spatial meaning, let $s\in\mathbb R^d$
and choose smooth real functions $e_1(s),\ldots,e_r(s)$.
For independent centered Gaussian variables $x_1,\ldots,x_r$ of variance
$h$, the random function
\[
 \phi(s)=\sum_{i=1}^r x_i e_i(s)
\]
is a finite-mode field. Each fixed spatial function $e_i$ is a mode,
and $x_i$ is its fluctuating amplitude. The Gaussian amplitudes make
this a free model. Passing to infinitely many modes requires precisely
the choices of coefficients and convergence just encountered.

\medskip

To attach positions to our polynomial operations, begin with translation.
Taylor's formula says
\[
 e^{z\partial_y}f(y)=f(y+z).
\]
Both sides are polynomials in $y,z$. Translation commutes with every
derivative, so it also respects $\star_h$.
Thus measurements can be translated and combined in a compatible way.
Singular coefficients appear when infinitely many modes participate.

On complex polynomials in $x_0,x_1,\ldots$, let $X_n$ multiply by $x_n$ and
let $D_n=\partial/\partial x_n$. Each polynomial uses finitely many
variables. The product rule gives
\[
 [D_m,X_n]=\delta_{mn}1,
 \qquad [A,B]:=AB-BA,
\]
where $\delta_{mn}$ is $1$ for $m=n$ and $0$ otherwise.
Package these operators into position-dependent series
\[
 A(z)=\sum_{n\ge0}D_nz^{-n-1},\qquad
 B(w)=\sum_{n\ge0}X_nw^n.
\]
Their products and commutator are interpreted coefficient by coefficient.
For each pair of exponents the operator is specified by one pair of
modes. Consequently
\[
 [A(z),B(w)]=\sum_{n\ge0}z^{-n-1}w^n\,1.
\]
The scalar coefficient is the geometric expansion of $1/(z-w)$
for $|w|<|z|$. An elementary differentiation-and-multiplication relation
has produced a pole at coincident positions.

The same rational function has another expansion:
\[
 C_>(z,w)=\sum_{n\ge0}z^{-n-1}w^n,
 \qquad
 C_<(z,w)=-\sum_{n\ge0}z^nw^{-n-1}.
\]
They converge respectively for $|w|<|z|$ and $|z|<|w|$.
Their difference makes sense as an array of coefficients of integer
powers of $z,w$:
\[
 \Delta(z,w)=C_>(z,w)-C_<(z,w)
 =\sum_{n\in\mathbb Z}z^{-n-1}w^n.
\]
Multiplication by $z-w$ cancels adjacent terms, giving
$(z-w)\Delta=0$ coefficientwise. Moreover, for every polynomial $f$,
\[
 [z^{-1}]\bigl(f(z)\Delta(z,w)\bigr)=f(w).
\]
Indeed, the term $az^j$ of $f$ selects exactly the term with $n=j$.
The symbol $[z^{-1}]$ means taking the coefficient of $z^{-1}$,
also called the residue coefficient. The difference between two expansion
orders thus evaluates a measurement at the collision.
The two convergent expansions have disjoint domains. Their coefficient
array comparison is a separate algebraic operation.

This example also shows why a single residue can lose information.
The Laurent polynomials $z^{-1}$ and $z^{-1}+z^{-2}$ have the same residue,
but multiplying them by $z$ gives different residues.
Keeping enough coefficients to permit the next multiplication matters
even when the eventual answer is one number.

Two complex position coordinates change the geometry of a collision.
Let $u,v$ be the differences of those coordinates. Distinct positions
satisfy $u\ne0$ or $v\ne0$.
On the first region allow finite Laurent sums in $u$ with polynomial
coefficients in $v$:
$R_u=\mathbb C[u,u^{-1},v]$.
On the second use $R_v=\mathbb C[u,v,v^{-1}]$.
Their intersection region admits
$R_{uv}=\mathbb C[u,u^{-1},v,v^{-1}]$.

A pair $a\in R_u$, $b\in R_v$ agrees on the intersection precisely
when $b-a=0$ there. A Laurent monomial lying in both rings has
nonnegative exponents in both variables. Thus every agreeing pair
comes from a polynomial in $u,v$.
There are nevertheless overlap terms that no difference $b-a$ produces:
\[
 \frac1{uv},\quad \frac1{u^2v},\quad\frac1{uv^2},\quad\ldots.
\]
Every monomial from $R_u+R_v$ has at least one nonnegative exponent.
Uniqueness of Laurent coefficients therefore gives the vector-space quotient
\[
 R_{uv}/(R_u+R_v)
 =\bigoplus_{a,b\ge0}\mathbb C\,u^{-a-1}v^{-b-1}.
\]
The direct sum means finite sums. These collision coefficients reside
in the comparison on the overlap, while agreeing functions themselves
extend across the origin.

Their pairing with polynomial measurements is explicit:
\[
 [u^{-1}v^{-1}]
 \bigl(u^iv^j u^{-a-1}v^{-b-1}\bigr)
 =\delta_{ia}\delta_{jb}
 \qquad(i,j,a,b\ge0).
\]
Multiplying a discarded term by a polynomial still leaves a nonnegative
exponent in at least one variable. Its $u^{-1}v^{-1}$ coefficient
therefore vanishes, so the pairing is well-defined on the quotient.
This explains how a singular coefficient can detect any prescribed
Taylor coefficient. There are two complex coordinates here, hence four
real coordinates. The overlap calculation describes a puncture in
$\mathbb C^2$ directly. Constructing free fields on this space also
requires a specified linear field equation and an inverse with fixed
normalization. These data determine the actual contraction coefficients.

\medskip

Spatial dependence places another condition on elimination: a
measurement has a region in which it is made.
A response $f$ has compact support if it vanishes outside some bounded
closed set $K$. For the finite-mode field $\phi$, a detector with
smooth compactly supported response $f(s)$ measures
\[
 M_f(\phi)=\int_{\mathbb R^d}f(s)\phi(s)\,ds.
\]
If $K$ lies in an open region $U$, the measurement depends only
on the field in $U$. This is the locality of this observable.
Its correlations follow from the Gaussian moments:
\[
 \begin{aligned}
 \langle M_fM_g\rangle
 &=h\sum_{i=1}^r
       \left(\int f(s)e_i(s)\,ds\right)
       \left(\int g(t)e_i(t)\,dt\right)\\
 &=\int\!\!\int f(s)C(s,t)g(t)\,ds\,dt,
 \qquad C(s,t)=h\sum_{i=1}^r e_i(s)e_i(t).
 \end{aligned}
\]
The function $C$ is the covariance kernel. Even detectors in disjoint
regions can be correlated: one mode may be nonzero in both regions.
Local dependence of a measurement and independence of its fluctuations
are different conditions.

Combining detector regions requires products as well as sums.
If $f=f_1+f_2$, with the two responses supported in respective regions,
then
\[
 M_f=M_{f_1}+M_{f_2},\qquad
 M_f^2=M_{f_1}^2+2M_{f_1}M_{f_2}+M_{f_2}^2.
\]
The middle term is a joint measurement using both regions.
Let $\mathcal O(U)$ be the algebra generated by $1$ and all $M_f$
with response supported compactly inside $U$.
If $U\subset V$, each such measurement also belongs to $\mathcal O(V)$.
For disjoint regions $U_1,\ldots,U_k\subset V$, set
\[
 \mu_{U_1,\ldots,U_k;V}(F_1,\ldots,F_k)=F_1\cdots F_k
 \in\mathcal O(V),\qquad F_i\in\mathcal O(U_i).
\]
Suppose $U_1,U_2\subset V_1$, $U_3\subset V_2$, and $V_1,V_2$
are disjoint subsets of $V$. Associativity says
\[
 \mu_{V_1,V_2;V}
 \bigl(\mu_{U_1,U_2;V_1}(F_1,F_2),F_3\bigr)
 =\mu_{U_1,U_2,U_3;V}(F_1,F_2,F_3).
\]
The same equality holds for every grouping. Permuting regions permutes
factors, and the empty product is $1$.
These maps and identities constitute a \emph{prefactorization algebra}.
A \emph{factorization algebra} additionally reconstructs observables
from regions and their overlaps, using covers that contain every finite
set of measurement positions in one member.
The finite model verifies the displayed product identities.
Reconstruction requires further overlap equations, including the mixed
measurements visible in $M_{f_1}M_{f_2}$.

\medskip

Symmetry imposes equations on these operations. Return to two real
Gaussian variables, now called $q,p$, with energy
$H=(q^2+p^2)/2$ and weight $e^{-H/h}$.
An infinitesimal rotation acts on polynomial observables by
\[
 Jf=q\frac{\partial f}{\partial p}
       -p\frac{\partial f}{\partial q}.
\]
It preserves the energy, since $JH=0$.
Integration by parts in $p$ and in $q$ gives
\[
 \langle q\,\partial_p f\rangle
   =h^{-1}\langle pqf\rangle
   =\langle p\,\partial_q f\rangle,
 \qquad \langle Jf\rangle=0.
\]
This is a symmetry identity for every observable, including products
of arbitrarily many measurements. Such identities are called Ward
identities. They express invariance after averaging.

The rotation has a Hamiltonian description. Write
$F_q=\partial F/\partial q$ and $F_p=\partial F/\partial p$.
For polynomials in $q,p$, define
\[
 \{F,G\}=F_qG_p-F_pG_q.
\]
Then $Jf=\{H,f\}$. More generally, $\{F,-\}$ is the infinitesimal
motion with velocity $(-F_p,F_q)$ in the $(q,p)$ plane, and
$\{F,F\}=0$ says that its own Hamiltonian is constant along that motion.
A monomial calculation gives a useful finite approximation. Fix an
integer $N\ge0$.
If every term of $F$ has degree at least two and every term of $g$
has degree at least $N+1$, then every term of $\{F,g\}$ has degree
at least $N+1$.
Thus these motions act on polynomials modulo terms of degree $N+1$
and higher. They fix the origin and act on the finite Taylor spaces.
A translation does not: $\partial_q(q^{N+1})=(N+1)q^N$ returns a
discarded term to the retained range.

This failure is measurable in the original multiplication and
differentiation operators. Let $V_N$ have basis $1,x,\ldots,x^N$.
Let $D$ differentiate, and let $X_N$ multiply by $x$ except that
$X_Nx^N=0$. For $0\le j<N$ the product rule still gives
\[
 (DX_N-X_ND)x^j=x^j.
\]
On the last basis vector it gives $-Nx^N$.
Let $\Pi_N$ keep only the term of degree $N$, so that
$\Pi_N(\sum_{j=0}^N a_jx^j)=a_Nx^N$. Then
\[
 [D,X_N]=1-(N+1)\Pi_N.
\]
The trace of a finite matrix is the sum of its diagonal entries.
Here it equals $(N+1)-(N+1)=0$.
This agrees with the general identity $\operatorname{tr}(AB)=
\operatorname{tr}(BA)$: expanding either trace gives the same finite
sum $\sum_{i,j}A_{ij}B_{ji}$.

On all polynomials, multiplication by $x$ instead gives $[D,X]=1$.
Every fixed polynomial eventually lies below the cutoff and sees this
identity exactly. The boundary term has disappeared from that
calculation, although its trace is $-(N+1)$.
The identity on all polynomials has no ordinary finite trace.
Hence agreement on each polynomial does not justify passing a trace
through the limit.
An infinite collection of matter modes requires a separate trace
construction before finite-dimensional symmetry calculations can
be carried over.

In a quantum theory, observable operators can fail to commute, as $D$
and $X$ do. A symmetry identity may acquire a correction when products
of such operators are defined. An \emph{anomaly} is a resulting defect
that the permitted changes of the theory cannot remove.
Our cutoff calculation only obstructs a trace argument.
To establish an anomaly for fields, one must specify their equations,
their averaging or operator rules, and the allowed corrections.
Local corrections must use fields and finitely many derivatives at
the same position. The existence of an unrestricted correction does
not ensure the existence of one with this property.

The force of the last requirement appears in elementary calculus.
For a smooth compactly supported function $f$ on the real line, put
\[
 F(t)=\int_{-\infty}^t f(s)\,ds.
\]
Then $F'=f$. To the left of the support, $F=0$.
To its right, $F=\int_{\mathbb R}f$.
Thus this primitive is compactly supported exactly when the total
integral of $f$ is zero. Even when it is compactly supported, its
value uses measurements away from $t$.

Could a formula using only finitely many derivatives at $t$ do the same
job for every $f$? Suppose such a linear formula were
\[
 Rf(t)=\sum_{j=0}^m a_j(t)f^{(j)}(t),\qquad (Rf)'=f,
\]
with fixed smooth coefficients $a_j$.
At any chosen point, compactly supported smooth functions can have
arbitrarily prescribed derivatives through order $m+1$.
Indeed, multiply a polynomial with those Taylor coefficients by a
smooth function equal to one near the point and zero outside a
larger interval.
Such functions can be constructed from $e^{-1/u}$ for $u>0$, extended
by zero for $u\le0$: every derivative tends to zero at the joining point.
In $(Rf)'$, the coefficient of $f^{(m+1)}$ is $a_m$.
It must vanish everywhere. Repeating the argument forces
$a_{m-1},\ldots,a_0$ to vanish, contradicting $(Rf)'=f$.
The integral gives a primitive, but no finite differential formula
provides a right inverse on all these test functions.

The spring identity led from a vanished boundary term to a product
correction. With fields, such corrections must also respect position,
support, and symmetry. The mathematical task is to construct them
together with all their composition equations, and to determine the
obstructions when they cannot exist.
